\section{Future Work}

SIDS 2026 \cite{katrakazas2026} announced a three-country pilot deployment of the BENEFITS framework across London, Tuscany, and Zaragoza, engaging approximately 150 participants accessing employment, healthcare, housing, and legal services. That pilot was announced against a system that, by SIDS 2026's own account, existed as specification rather than code. This paper is what had to happen between the announcement and the deployment: the mechanisms the pilot will run on now exist as a built reference implementation, and six of them have been measured under controlled conditions rather than assumed correct from their design documents. That makes this paper the technical readiness gate for the pilot, not a substitute for it: the results in Chapter 5 hold in simulation, and do not yet say whether they survive three real deployments, three legal frameworks, and populations who did not co-design the system now mediating their data. Five open threads follow, each traceable to a specific finding here or an unresolved question SIDS 2026 itself raised.

\subsection{From Simulation to Field Conditions}

Every quantitative result in Chapter 5 comes from a seeded simulation harness. E1's and E3's timings were taken inside a container throttled to one CPU core and 1 GB of memory, chosen to approximate a low-end handset rather than measured on the device population of three real sites. E2's zero-leaked-record result assumed an online-probability parameter standing in for whatever connectivity pattern the sites actually exhibit. E3's proof-size and verification-time curves assume an epoch-sealing cadence chosen for the experiment, not tuned to real contribution rates. E4's and G2's privacy-utility tradeoffs used synthetically generated populations at fixed sizes (30, 100, 200, 300) that may not match the per-site cohorts the pilot recruits. And G1's quorum-availability comparison assumed steward action probabilities chosen to separate two policies cleanly in a sweep, not measured from stewards with real jobs and real availability.

None of this invalidates the mechanism-level findings, but the pilot needs its own instrumentation before any of these numbers can be claimed to hold in the field. Concretely: log real device online/offline duty cycles per site to check the zero-leak result against whatever connectivity floor the least-connected site exhibits. Log real population sizes and consent-heterogeneity levels at each query the privacy pipeline serves, to locate each site on the E4 and G2 curves rather than assume it falls inside the simulated range. Log real steward response latency and action frequency, to check G1's quorum-failure predictions against actual behavior. And measure signature-verification and proof-generation time on the actual devices participants own and inside the production consent filter and audit path, rather than the throttled profile and separate harness E1 and E3 used as a stand-in for one. Where field numbers diverge from simulated ones, that divergence is a finding, not a correction to fold away quietly.

\subsection{Closing the Steward-Quorum Gaps}

Chapter 3 documented, and Chapter 6 restates as a limitation, that the steward-quorum release gate as built carries two separate gaps against what SIDS 2026 specified. The first is cryptographic: \texttt{StewardQuorumFunction} counts distinct steward signatures against an ODRL constraint threshold rather than gating release on reconstructing a decryption key from steward-held shares, so the guarantee depends on trusting the connector operator to enforce the check honestly rather than on cryptographic impossibility independent of that operator. The second is the veto: SIDS 2026's own text specifies that any single steward may veto a release, returning it to the requester with a written justification, and \texttt{StewardQuorumFunction.evaluate} never reads the recorded veto list at all, so a vetoed release with a satisfied signature count clears contract negotiation exactly as an unvetoed one would. Both are real gaps between what was claimed and what was built, and neither is resolved by this paper. Each is future work, with a different character.

The veto gap is the more urgent of the two to close, not because it costs more, G1 (section 5.5) shows precisely how much enforcing it would cost, but because it diverges from a guarantee SIDS 2026 already published to its readers rather than from an internal design intent a reference implementation stopped short of realizing, and because closing it is a comparatively small engineering task: \texttt{StewardAuthorization} already carries the \texttt{vetoes} list \texttt{StewardQuorumFunction.evaluate} currently ignores, so the fix is adding a check of that list to the existing constraint function, not building new infrastructure. Doing so would trade away exactly the availability G1 measured Arm B losing relative to Arm A at every tested quorum fraction and action probability, a real cost this paper's contribution is to have made visible before the pilot has to absorb it silently.

The cryptographic gap requires more substantial work. Closing it means implementing genuine threshold decryption over the release artifact, following one of the two closest technical comparators already identified in Chapter 3: CALYPSO's on-chain secret management, where decryption requires an actual threshold of a trustee committee's participation, benchmarked directly against committee size \cite{kokoriskogias2020}. Or MedCo's federated key-share model, in which decrypting a pooled result requires the collective participation of every site holding a share, deployed for real clinical and genomic data sharing \cite{raisaro2019}. Either pattern would require reworking \texttt{ReleaseArtifact} to hold ciphertext under a quorum-reconstructable key rather than plaintext behind a policy gate, distributing key shares to stewards through a channel the current \texttt{StewardAuthorization} record does not need, and handling share rotation when a steward leaves. The corresponding experiment this would enable, and that the current architecture cannot support, is a direct latency and availability comparison between the built policy-check gate and a threshold-decryption gate under the same quorum-fraction and action-probability sweep already run for G1, extended to include a veto-enforcing variant of each once the veto gap above is closed, showing concretely what cryptographic independence from the operator costs in practice.

\subsection{From Epsilon to Deliberation}

SIDS 2026's Open Governance Challenges section raised a concern it left explicitly unresolved: at epsilon = 0.5 per query and a total epsilon budget of 2.0 per site, differential privacy provides a formal guarantee but can render aggregate statistics analytically meaningless at the population sizes individual service sites actually have, undermining the political value that motivated participation. SIDS 2026 proposed adaptive privacy budgets that involve community members directly in the tradeoff, translating abstract epsilon values into concrete scenario descriptions of what an analyst can and cannot learn, rather than asking a steward to reason about epsilon as a number.

E4's surviving-utility curves are a first quantitative answer to what those scenario descriptions would need to say. At population = 30, the smallest site size tested, surviving median relative error at epsilon = 0.1 ranges from 0.513 to 1.127 across k thresholds, and suppression removes 64.7 to 95.6 percent of query attempts at that same epsilon. At population = 300, the same epsilon leaves substantially more usable, with suppression as low as 6.5 percent at k = 3. That is the concrete tradeoff SIDS 2026's proposed scenario descriptions would need to render legible: not "epsilon is 0.1" but "at this site's population, roughly two in three queries at this budget will return nothing at all." The direct follow-on is building that translation layer, presenting stewards with the suppression rate and expected error at candidate budgets for their own site's actual population, and running the adaptive-budget deliberation SIDS 2026 proposed as a real decision process in the pilot rather than a fixed epsilon set centrally. An empirical question this paper's simulated curves inform but cannot answer is whether stewards given this information choose materially different budgets and whether communities judge the resulting aggregates more useful.

\subsection{Sustainable Steward Governance Across Three Sites}

SIDS 2026 identified a structural tension in the steward role: independence from service providers and researchers gives a steward's veto its governance value, but the financial support a steward needs to sustain the role risks creating dependency on the organizations they are meant to hold accountable. SIDS 2026 named three candidate models to compare, a stipend model (200 euros per month for 10 hours), in-kind support, and institutional embedding as part-time staff, with the comparison deferred to the three-country pilot.

G1 bears on that comparison without resolving it. G1 does not show that a stricter hypothetical policy costs more than the real one. It shows what enforcing the veto guarantee SIDS 2026 already specified, currently unenforced in the built system, would cost in quorum availability, with a minimum divergence of 0.555 at the narrowest panel between the built, veto-blind policy and the veto-enforcing one SIDS 2026's own text describes. The p\_act parameter in that sweep stands in for how available and responsive a steward actually is, which is exactly what the three governance models would be expected to affect differently, and exactly what determines how expensive closing the veto gap would be at a given site: a well-resourced, responsive steward pool absorbs an enforced veto's availability cost differently than an under-resourced one does. Even the built system's current, veto-blind policy is not uniformly cheap on its own terms: at quorum\_fraction = 0.7, its failure rate already ranges from 0.340 to 0.440 across the swept action-probability range, from stewards simply not signing in time, before any veto enforcement or any real-world reduction in steward availability from an under-resourced governance model is accounted for. Comparing the three governance models empirically across London, Tuscany, and Zaragoza should therefore include measuring each site's actual steward action rate against the quorum-availability model already built here, once the veto gap described in Section 8.2 is closed, so that whichever model a site adopts is evaluated against the guarantee SIDS 2026 actually specified rather than against the one the current implementation happens to enforce. The independence-versus-dependency tradeoff itself is a governance question outside this paper's scope. The availability consequence of whichever model a site adopts, under the veto guarantee actually specified, is not, and the instrumentation to check it already exists in G1's harness.

\subsection{From Sovereignty Mechanisms to Service Decisions}

SIDS 2026 posed three research questions for the pilot. Do steward-sovereignty controls designed with Greek migrants transfer across legal frameworks and cultural norms. Do observable sovereignty mechanisms replicate the trust and participation effects measured with participants who did not co-design the system. Does community-governed participatory data actually influence service decisions, measured against control sites that rely only on traditional administrative records. The third question is the one the pilot was designed to answer and the one this paper cannot touch. Nothing in Chapter 5 speaks to whether a service provider changes a decision because the data behind it passed through a community-governed pipeline rather than an administrative one.

What this paper changes is what stands underneath that question when the pilot finally asks it. SIDS 2026's own evaluation was necessarily a trust measurement, participants reporting confidence in mechanisms that existed on paper. A finding that community-governed data changed service decisions would previously have rested on an architecture whose guarantees were asserted rather than measured. It now rests on one where signature-verification cost (measured in a harness standing in for the production filter), revocation propagation under partition, audit-chain proof size and verification cost, privacy-utility tradeoffs, quorum availability, and consent-heterogeneity effects have each been characterized directly, including where the architecture falls short of its own original design, as with the steward-quorum gaps disclosed in Section 8.2. Tamper detection itself was established separately, at the mechanism level, by an earlier implementation phase's unit tests, not by this paper's own experiments. That does not make the outcome-influence question easier to answer, and it does not by itself confound-proof the pilot's result. This paper closes the confound of whether the specified mechanisms work at all, under any condition. But it closes that confound only under simulated conditions: section 8.1 above and Chapter 6's own Limitations state plainly that every number behind the claim comes from simulation, not field measurement. What remains open, and is exactly what section 8.1 says the pilot's own instrumentation must check, is whether these mechanisms hold under real field conditions, three real deployments, three legal frameworks, and populations who did not co-design the system. This paper closes the first confound. The second is the pilot's to close.
